% sec3_9.tex — Section 3.9: Application: Returns from Schooling

Since Mincer's (1958) pioneering study, the relationship between the wage rate
and schooling has been the subject of a large number of empirical and theoretical
investigations.  You might find the amount of attention puzzling because the
explanation of the positive relationship seems to be obvious: education enhances
the individual's productivity.  There are, however, other explanations.  For example,
in the Spencian job market signaling model, more educated individuals receive
higher wages only because education is used as a signal of higher ability; although
education does not increase the individual's earning capacity, there is a correlation
between the wage rate and schooling because both variables are influenced by the
third variable, ability.  This section shows how to use the technique of this chapter
to isolate the effect of education on the wage rate from that of ability.  One of the
earliest studies to address this issue is Griliches (1976).

%% ─── The NLS-Y Data ───────────────────────────────────────────────────────
\subsection*{The NLS-Y Data}

The data used by Griliches are the Young Men's Cohort of the National Longitudinal
Survey (NLS-Y).  This cohort was first surveyed in 1966 at ages 14--24, with
5,225 respondents, and was resurveyed at one- or two-year intervals thereafter.  By
1969, about a quarter of the original sample was lost, but there are 2,026 individuals
who reported earnings in 1969 and whose records are complete enough to
allow derivation of all the variables for analysis.  A very attractive feature of the
NLS-Y is its inclusion of two measures of ability.  One of them is the score on
the Knowledge of the World of Work (KWW) test administered by the NLS interviewers
in 1966.  The other measure is the IQ score.  All youths in the survey who
had completed ninth grade by 1966 were asked to sign waivers letting their school
supply the survey administrator their scores on various tests and other background
materials.  The resulting School Survey, conducted in 1968, yielded data on different
mental ability scores, which were combined into IQ equivalents.  Of the 2,026
individuals with information about the 1969 wage rate and other variables, the IQ
score is available for 1,362 individuals, reflecting the fact that the School Survey
was able to cover only two-thirds of the original sample.  Our discussion here concerns
Griliches's results based on this smaller sample.  Table~3.1 reports the means
and standard deviation of key variables.

%% ─── Table 3.1 ─────────────────────────────────────────────────────────────
\begin{table}[htbp]
\centering
\caption{Characteristics of Young Men from the National Longitudinal Survey}
\label{tab:3.1}
\medskip
Means and Standard Deviations (in parentheses)
\medskip

\begin{tabular}{@{}l r@{}}
\toprule
Variable & \\
\midrule
Sample size & 1{,}362 \\[4pt]
Age in 1969 & 22.3 \\
 & (3.2) \\[4pt]
Schooling in years in 1969 ($S$) & 12.5 \\
 & (1.9) \\[4pt]
Logarithm of hourly wages (in cents) in 1969 ($LW$) & 5.68 \\
 & (0.40) \\[4pt]
Score on the Knowledge of the World of Work test ($KWW$) & 35.1 \\
 & (7.9) \\[4pt]
IQ score ($IQ$) & 97.7 \\
 & (15.3) \\[4pt]
Experience in years in 1969 & 3.7 \\
 & (2.8) \\
\bottomrule
\end{tabular}

\smallskip
{\footnotesize Source: Griliches (1976, Table~1).}
\end{table}

%% ─── The Semi-Log Wage Equation ────────────────────────────────────────────
\subsection*{The Semi-Log Wage Equation}

The typical wage equation estimated in the literature is the semi-log form:
\begin{equation}
  LW = \alpha + \beta S + \gamma A + \boldsymbol{\delta}'\mathbf{h} + \varepsilon,
  \label{eq:3.9.1}
\end{equation}
where $LW$ is the log wage rate for the individual, $S$ is schooling in years, $A$ is a
measure of ability, $\mathbf{h}$ is the vector of observable characteristics of the individual
(such as experience and location dummies), $\boldsymbol{\delta}$ is the associated vector of coefficients,
and $\varepsilon$ is the unobservable error term with zero mean (in this section, the
individual subscript $i$ will be dropped for notational simplicity).  The semi-log
specification for schooling $S$ is often justified by appealing to the well-established
stylized fact from large cross-section data (such as the Current Population Survey)
that the relationship between log wages and schooling is linear.\footnote{For the
functional form issue, see Card (1995).}

The schooling coefficient $\beta$ measures the percentage increase in the wage rate
the individual would receive if she had one more year of education.  It therefore
represents the marginal return from investing in human capital, which should be in
the same order of magnitude as the rate of return from financial assets.  We assume
that the nonconstant regressors, $(S, A, \mathbf{h})$, are uncorrelated with the error $\varepsilon$
so that the OLS is an appropriate estimation technique if ability is included in
the regression along with $S$ and $\mathbf{h}$.  In the rest of this section, we examine what
Griliches called the ``ability bias'' --- biases on the OLS estimate of $\beta$ that would
arise when ability $A$ is not included in the regression and when its imperfect measure
is included in its place.

%% ─── Omitted Variable Bias ─────────────────────────────────────────────────
\subsection*{Omitted Variable Bias}

Sometimes the data set you work with has no measures of $A$ (this is true for the
Current Population Survey, for example).  What is the consequence of ignoring
ability by omitting $A$ from the wage equation?  We know from Section~2.9 that
the regression of $LW$ on a constant, $S$, and $\mathbf{h}$ provides a consistent estimator of the
corresponding least squares projection, which can be written as
\begin{align}
  \tilde{\E}^*(LW \mid 1, S, \mathbf{h})
  &= \tilde{\E}^*(\alpha + \beta S + \gamma A + \boldsymbol{\delta}'\mathbf{h}
    + \varepsilon \mid 1, S, \mathbf{h})
  \quad \text{(by \eqref{eq:3.9.1})} \notag \\
  &= \alpha + \beta S + \boldsymbol{\delta}'\mathbf{h}
    + \gamma\, \tilde{\E}^*(A \mid 1, S, \mathbf{h})
    + \tilde{\E}^*(\varepsilon \mid 1, S, \mathbf{h}).
  \label{eq:3.9.2}
\end{align}
Since the regressors in \eqref{eq:3.9.1} are all predetermined, we have $\E(\varepsilon) = 0$,
$\E(S\,\varepsilon) = 0$, and $\E(\mathbf{h}\cdot\varepsilon) = \mathbf{0}$.  So
$\tilde{\E}^*(\varepsilon \mid 1, S, \mathbf{h}) = 0$.  Writing
$\tilde{\E}^*(A \mid 1, S, \mathbf{h}) = \theta_1 + \theta_S\, S + \boldsymbol{\theta}_{\mathbf{h}}'\mathbf{h}$,
\eqref{eq:3.9.2} becomes
\[
  \tilde{\E}^*(LW \mid 1, S, \mathbf{h})
  = (\alpha + \gamma \theta_1) + (\beta + \gamma \theta_S)\, S
  + (\boldsymbol{\delta} + \gamma \cdot \boldsymbol{\theta}_{\mathbf{h}})'\mathbf{h}.
\]
Therefore, the OLS coefficient estimate can be asymptotically biased for \emph{all} the
included regressors.  This phenomenon is called the \textbf{omitted variable bias}.  In
particular,
\[
  \plim\, \hat{\beta}_{\text{OLS}} = \beta + \gamma \theta_S.
\]
That is, the OLS estimate of the schooling coefficient $\beta$ includes the indirect effect
of ability through schooling ($\gamma \theta_S$) as well as the direct effect of schooling
($\beta$).  If $\theta_S$ is positive, then $\hat{\beta}_{\text{OLS}}$ is asymptotically biased upward.

Using the sample of 1,362 individuals from the NLS-Y described above,
Griliches estimated the wage equation for 1969.  The list of variables included in
$\mathbf{h}$ will not be given here; suffice to say that it includes experience in years and some
region and city-size dummies.  Griliches's estimate of the schooling coefficient
when ability is ignored in the wage equation is reproduced on line~1 of Table~3.2.
On the face of it, the estimate looks good: the point estimate resembles the rate of
return one might get from financial assets, and it is sharply estimated as evidenced
by the high $t$-value.

%% ─── IQ as the Measure of Ability ──────────────────────────────────────────
\subsection*{IQ as the Measure of Ability}

As already mentioned, the NLS-Y has two measures of ability, $KWW$ (the score on
the Knowledge of the World of Work test) collected in 1966 and $IQ$ (the IQ score).
Since the NLS-Y respondents were at least fourteen years old in 1966, $KWW$ would
reflect the effect of schooling already undertaken, and so it cannot be a measure of
raw ability.  The IQ score does not have this problem.  If $IQ$ were a perfect measure
of ability so that $A$ can be equated with it, then the wage equation \eqref{eq:3.9.1}
could be estimated consistently with $IQ$ substituting for $A$.  Griliches's OLS estimates
of $\beta$ (schooling coefficient) and $\gamma$ (ability coefficient) when $IQ$ is included in
the equation are reported in line~2 of Table~3.2.  Now the estimated schooling coefficient
is lower, confirming our prediction that the estimated schooling coefficient includes
the effect of ability on the wage rate when ability is omitted from the regression.

%% ─── Errors-in-Variables ──────────────────────────────────────────────────
\subsection*{Errors-in-Variables}

Of course the IQ score may not be an error-free measure of ability.  If $\eta$ is the
measurement error, $IQ$ is related to $A$ as
\begin{equation}
  IQ = \phi + A + \eta,
  \label{eq:3.9.3}
\end{equation}
with $\E(\eta) = 0$.  The interpretation of $\eta$ as measurement error makes it reasonable
to assume that $\eta$ is uncorrelated with $A$, $S$, $\mathbf{h}$, and the wage equation error term
$\varepsilon$.  Substituting \eqref{eq:3.9.3} into \eqref{eq:3.9.1}, we obtain
\begin{equation}
  LW = (\alpha - \gamma\phi) + \beta S + \gamma\, IQ + \boldsymbol{\delta}'\mathbf{h}
  + (\varepsilon - \gamma\eta).
  \label{eq:3.9.4}
\end{equation}

We now illustrate for this example the general result that, if at least one of the
regressors is measured with error, the OLS estimates of \emph{all} the regression coefficients
can be asymptotically biased.

To examine the consequence of using the error-ridden measure $IQ$ for $A$, consider
the corresponding least squares projection:
\begin{align}
  \tilde{\E}^*(LW \mid 1, S, IQ, \mathbf{h})
  &= (\alpha - \gamma\phi) + \beta S + \gamma\, IQ + \boldsymbol{\delta}'\mathbf{h} \notag \\
  &\quad + \tilde{\E}^*(\varepsilon \mid 1, S, IQ, \mathbf{h})
    - \gamma\, \tilde{\E}^*(\eta \mid 1, S, IQ, \mathbf{h}).
  \label{eq:3.9.5}
\end{align}
Now consider $\tilde{\E}^*(\varepsilon \mid 1, S, IQ, \mathbf{h})$ in \eqref{eq:3.9.5}.
Since $\E(\varepsilon) = 0$ and $(S, \mathbf{h})$ are uncorrelated with $\varepsilon$,
$(1, S, \mathbf{h})$ are orthogonal to $\varepsilon$.  $IQ$ is also orthogonal to $\varepsilon$ because
\[
  \E(IQ \cdot \varepsilon) = \E[(\phi + A + \eta)\,\varepsilon]
  \quad \text{(by \eqref{eq:3.9.3})}
  = \E(\eta\,\varepsilon) = 0.
\]
Thus, $\tilde{\E}^*(\varepsilon \mid 1, S, IQ, \mathbf{h}) = 0$ in \eqref{eq:3.9.5}.

So the biases, if any, equal
$-\gamma\, \tilde{\E}^*(\eta \mid 1, S, IQ, \mathbf{h})$ in \eqref{eq:3.9.5}.
Let $(\theta_S, \theta_{IQ}, \boldsymbol{\theta}_{\mathbf{h}})$ be the projection coefficients of
$(S, IQ, \mathbf{h})$ in $\tilde{\E}^*(\eta \mid 1, S, IQ, \mathbf{h})$.  By the formula
(2.9.7) from Chapter~2,
\begin{equation}
  \begin{bmatrix} \theta_S \\ \theta_{IQ} \\ \boldsymbol{\theta}_{\mathbf{h}} \end{bmatrix}
  = \begin{bmatrix}
      \Var(S) & \Cov(S, IQ) & \Cov(S, \mathbf{h}') \\
      \Cov(IQ, S) & \Var(IQ) & \Cov(IQ, \mathbf{h}') \\
      \Cov(\mathbf{h}, S) & \Cov(\mathbf{h}, IQ) & \Var(\mathbf{h})
    \end{bmatrix}^{-1}
  \begin{bmatrix} \Cov(S, \eta) \\ \Cov(IQ, \eta) \\ \Cov(\mathbf{h}, \eta)
  \end{bmatrix}.
  \label{eq:3.9.6}
\end{equation}
In this expression, $\Cov(S, \eta)$ and $\Cov(\mathbf{h}, \eta)$ are zero by assumption.
$\Cov(IQ, \eta)$, however, is not zero because
\begin{align*}
  \Cov(IQ,\, \eta) &= \E(IQ\;\eta) \quad \text{(since $\E(\eta) = 0$)} \\
  &= \E[(\phi + A + \eta)\,\eta] \quad \text{(by \eqref{eq:3.9.3})} \\
  &= \phi\,\E(\eta) + \E(A\,\eta) + \E(\eta^2) \\
  &= \Var(\eta) \quad \text{(since $\E(\eta) = 0$ and $\Cov(A,\eta) = 0$)}.
\end{align*}
That is, the measurement error, if uncorrelated with the true value, is necessarily
correlated with the measured value.  Using the fact that
\[
  \Cov(S, \eta) = 0, \quad \Cov(IQ, \eta) = \Var(\eta), \quad
  \Cov(\mathbf{h}, \eta) = \mathbf{0},
\]
the projection coefficients can be rewritten as
\[
  \begin{bmatrix} \theta_S \\ \theta_{IQ} \\ \boldsymbol{\theta}_{\mathbf{h}} \end{bmatrix}
  = \Var(\eta) \cdot \mathbf{a},
\]
where
\[
  \mathbf{a} \equiv \text{second column of }
  \begin{bmatrix}
    \Var(S) & \Cov(S, IQ) & \Cov(S, \mathbf{h}') \\
    \Cov(IQ, S) & \Var(IQ) & \Cov(IQ, \mathbf{h}') \\
    \Cov(\mathbf{h}, S) & \Cov(\mathbf{h}, IQ) & \Var(\mathbf{h})
  \end{bmatrix}^{-1}.
\]

Since the second element of $\mathbf{a}$ is positive (it is a diagonal element of the inverse of
a variance-covariance matrix), the OLS estimate of the ability coefficient is biased
downward.  The direction of the asymptotic bias for the schooling coefficient, however,
depends on the sign of the first element of $\mathbf{a}$.  Typically, $\Cov(S, IQ)$ would
be positive, so, barring unusually strong correlation of $(S, IQ)$ with $\mathbf{h}$, the first
element of $\mathbf{a}$ would be negative.  Thus, we conjecture that $\hat{\beta}_{\text{OLS}}$ is
biased upward for schooling.

Therefore, in the regression of $LW$ on a constant, $S$, $IQ$, and $\mathbf{h}$,
\begin{align}
  \plim\, \hat{\beta}_{\text{OLS}}
  &= \beta - \gamma \cdot \Var(\eta) \cdot \text{(1st element of $\mathbf{a}$)},
  \label{eq:3.9.7a} \\
  \plim\, \hat{\gamma}_{\text{OLS}}
  &= \gamma - \gamma \cdot \Var(\eta) \cdot \text{(2nd element of $\mathbf{a}$)}.
  \label{eq:3.9.7b}
\end{align}

%% ─── 2SLS to Correct for the Bias ─────────────────────────────────────────
\subsection*{2SLS to Correct for the Bias}

To control for the bias, Griliches applies the 2SLS to the wage equation.  To do
so, the set of instruments needs to be specified.  The predetermined regressors,
$(1, S, \mathbf{h})$, can be included in the set.  The additional variables included in the set
to instrument $IQ$ are age, age squared, $KWW$, mother's education, and some other
background variables of the individual (such as father's occupation).  Those variables
are thus assumed to be predetermined.  Griliches's 2SLS estimate for this
specification is reproduced in line~3 of Table~3.2.\footnote{The $R^2$ is not reported for
the 2SLS estimate because, unlike in the OLS case, the sum of squared 2SLS
residuals cannot be divided between the ``explained variation'' and the ``unexplained
variation.''}  In accordance with our prediction that the OLS estimate of the ability
coefficient is biased downward, the OLS estimate of the ability coefficient of
0.0019 in line~2 is lower than the 2SLS estimate of 0.0038 in line~3.  Our conjecture
that $\hat{\beta}_{\text{OLS}}$ of the schooling coefficient when $IQ$ is included in the
regression is biased upward, too, is borne out by data because the estimate of
$\beta$ of 0.059 in line~2 is higher than the estimate of 0.052 in line~3.

%% ─── Table 3.2 ─────────────────────────────────────────────────────────────
\begin{table}[htbp]
\centering
\caption{Parameter Estimates}
\label{tab:3.2}
\medskip
Estimation equation: $LW = \alpha + \beta S + \gamma\, IQ + \text{other controls}$
\medskip

\begin{tabular}{@{}c l cc cc l l@{}}
\toprule
& Estimation & \multicolumn{2}{c}{Coefficient of}
& & & Which regressor & Excluded \\
\cmidrule(lr){3-4}
Line no.\ & technique & $S$ & $IQ$ & $SER$ & $R^2$
& is endogenous? & predetermined variables \\
\midrule
1 & OLS & 0.065 & --- & 0.332 & 0.309 & none & --- \\
  &     & (13.2) & & & & & \\[4pt]
2 & OLS & 0.059 & 0.0019 & 0.331 & 0.313 & none & --- \\
  &     & (10.7) & (2.8) & & & & \\[4pt]
3 & 2SLS & 0.052 & 0.0038 & 0.332 & --- & $IQ$ & $MED$, $KWW$, age, \\
  &      & (7.0) & (2.4) & & & & age squared, back- \\
  &      &       &        & & & & ground variables \\
\bottomrule
\end{tabular}

\smallskip
{\footnotesize Source: Line~1: equation~(B1) in Griliches (1976, Table~2).
Line~2: equation~(B3) in Griliches's Table~2.
Line~3: line~3 in Griliches's Table~5.
Figures in parentheses are $t$-values rather than standard errors.}
\end{table}

To summarize, the ``ability bias'' studied by Griliches is that the schooling
coefficient is biased upward if ability is ignored and is biased in an unknown direction
(but perhaps upward) if an imperfect measure of ability ($IQ$ in the present case)
is included.  The 2SLS provides a solution, but it is predicated on the assumption
that the instruments used (such as $MED$ and $KWW$) are uncorrelated with unobserved
wage determinants $\varepsilon$ and $\eta$.  Those determinants would include personal
characteristics like diligence.  It seems not unreasonable to suppose that $KWW$
cannot serve as an instrument.  If so, $KWW$ is not predetermined either, because some
of the mother's personal characteristics that influenced her education and that are
valued in the marketplace would be inherited by the child.  Whether those variables
included in the set of instruments are predetermined can be tested by Sargan's
statistic.  But the statistic is not reported in Griliches's paper (which was written
long before doing so became a standard practice).

%% ─── Subsequent Developments ──────────────────────────────────────────────
\subsection*{Subsequent Developments}

Perhaps because it is so difficult to come up with valid instruments that are uncorrelated
with unobserved characteristics but correlated with $IQ$, controlling for the
``ability bias'' continued to attract much attention in the literature.  There is a fairly
large literature starting with Behrman et al.\ (1980), which compares identical twins
with different levels of education to control for genetic characteristics and family
background.  One can also compare the same individual at two points in time.  In
either case, the appropriate estimation technique is what is called the fixed-effect
estimator, to be covered in Chapter~5.

The endogeneity of schooling is another major issue.  If one takes the view that
the error term includes a host of unobservable individual characteristics that might
affect the individual's choice of schooling, then schooling needs to be treated as an
endogenous variable.  But again, finding valid instruments uncorrelated with unobservable
characteristics but correlated with schooling is extremely difficult.  The
recent literature can be viewed as a search for the determinants of schooling having
little to do with the individual's unobserved characteristics.  The variable used
in Angrist and Krueger (1991), for example, is whether the individual is subject to
compulsory schooling laws.

Finally, the literature on the relationship between school quality and earnings
is attracting renewed attention.  See Card and Krueger (1996) for a survey.

%% ─── Questions for Review ──────────────────────────────────────────────────
\bigskip
\begin{center}\rule{0.6\textwidth}{0.4pt}\end{center}
\noindent\textsc{Questions for Review}
\medskip

\begin{enumerate}
\item List all the assumptions made in this section about the means and covariances
  of $(S, A, IQ, KWW, \mathbf{h}, \varepsilon, \eta)$.  Study how Griliches states those assumptions.
  Which ones are crucial for the 2SLS to be consistent?

\item Suppose the relationship between $IQ$ and $A$ is given not by \eqref{eq:3.9.3} but by
  \[
    IQ = \phi + \pi \cdot A + \eta.
  \]
  Is the 2SLS estimator of $\beta$ (the schooling coefficient) still consistent?
  [Answer: Yes.]
\end{enumerate}
